\documentclass[10pt,a4paper]{amsart}
\usepackage[margin=19mm]{geometry}
\usepackage{amsmath,amssymb,mathtools}
\usepackage[scr=rsfs,cal=euler]{mathalfa}
\usepackage{xfrac}
\usepackage[unicode,hidelinks,bookmarks=false]{hyperref}
\newcommand{\ie}{i.e.\ }
\newcommand{\cf}{cf.\ }
\newcommand{\resp}{respectively\ }
\newcommand{\Subsection}{\subsection}
\providecommand{\nobreakdash}{\nobreak\hbox{-}}
\setlength{\emergencystretch}{2em}
\multlinegap0pt
\usepackage[shortlabels]{enumitem}
\usepackage{amssymb}
\usepackage{tikz}
\usepackage{dsfont}
\newtheorem{lemma}{Lemma}
 \newcommand{\cat}[1]{\ensuremath{\mathsf{#1}}} % Category
\def\cl{\mathfrak{c}}
\def\op{\mathfrak{o}}
\newcommand{\tbigvee}{\mathop{\textstyle \bigvee }}%%
\newcommand{\tbigwedge}{\mathop{\textstyle \bigwedge }}%%
\title{The lattice of smooth sublocales as a Bruns--Lakser completion}
\author{Igor Arrieta, Anna Laura Suarez}
\date{Source: arXiv:2602.15583v3 (CC BY 4.0)}
\begin{document}
\maketitle

\noindent\emph{Source and licence.} The lattice of smooth sublocales as a Bruns--Lakser completion. Version arXiv:2602.15583v3 (CC BY 4.0), distributed by arXiv under the Creative Commons Attribution 4.0 International licence, http://creativecommons.org/licenses/by/4.0/, as recorded in the arXiv licence metadata for this exact version and shown on the abstract page https://arxiv.org/abs/2602.15583v3. Full licence notice: \texttt{licenses/arxiv-2602.15583-CC-BY-4.0.txt}.\par\smallskip

\noindent\textit{Editorial scope.} Counted unit: the article's lemma phipsi (source0.tex lines 684--686). The article's complete original proof of it is delivered with it (source0.tex lines 688--726, one \texttt{proof} environment of 39 lines). No mathematics is written, repaired or re-derived: the statement and the proof are the article's own lines, in the article's own notation, and no retained line is substituted or re-flowed.  Retained uncounted as the source-local context the proof needs: lines 295--311.  Nothing was omitted from any retained span, so the package carries no omission record.  No retained line contains a citation, so the wrapper carries no bibliography and no bibliography entry.  No retained line contains a figure environment, an image-inclusion command or drawing code, so no image is delivered and the package declares no asset dependency.  Editorial formatting only, in addition: an AMS \texttt{amsart} preamble that restates the article's own declarations and macros used by the retained lines, namely the environments \texttt{lemma} on separate counters, together with the standard packages those lines need.  The retained environments print in this document as Lemma 1 in order of appearance, whereas the article prints the same items with the numbering of its own full document.  The complete native source of the article as distributed in the arXiv e-print for this exact version is bundled unchanged as \texttt{sources/194808/source0.tex}.
\subsection{Some Heyting rules}
\addtocontents{toc}{\protect\setcounter{tocdepth}{-1}}
For the reader's convenience, we list here some of the properties satisfied by the Heyting operator in a frame $L$. For any $a,b,c\in L$ and any $\{a_i\}_{i\in I}\subseteq L$, the following hold:
\begin{enumerate}[label=\textup{(H\arabic*)},leftmargin=2.0\parindent]
\item \label{H1} $1\to a=a$\textup;
\item \label{H2} $a\leq b$ if and only if $a\to b=1$\textup;
\item \label{H3} $a\leq b\to a$\textup;
\item \label{H4} $a\to b=a\to (a\wedge b)$\textup;
\item \label{H5} $a\wedge (a\to b)=a\wedge b$\textup;
\item \label{H6} $a\wedge b=a\wedge c$ if and only if $a\to b=a\to c$\textup;
\item \label{H7} $(a\wedge b)\to c=a\to (b\to c)= b\to (a\to c)$\textup;
\item \label{H8} $a=(a\vee b)\wedge (b\to a)$\textup;
\item \label{H9} $a\leq (a\to b)\to b$\textup;
\item \label{H10} $((a\to b)\to b)\to b=a\to b$\textup;
\item \label{H11}$(\tbigvee_{i\in I}a_i)\to b= \tbigwedge_{i\in I}(a_i\to b)$\textup;
\item \label{H12}$b\to (\tbigwedge_{i\in I}a_i)= \tbigwedge_{i\in I}(b\to a_i)$\textup.
\end{enumerate}

\begin{lemma}\label{l: fixpoints phipsi}
The relation $\varphi(\psi(U))=U$ holds if and only if $U$ is admissible.
\end{lemma}

\begin{proof}
Assume first that  $\varphi(\psi(U))=U$, then $U=\{ (a,b)\in \mathsf{LC}(L)\mid \cl(a)\cap\op(b)\subseteq \tbigvee_{(c,d)\in U}\cl(c)\cap\op(d)\}$. Now if $\{(a_i,b_i)\}\subseteq U$ is a locally exact family, then $\tbigvee_{i\in I}\cl(a_i)\cap\op(b_i)$ is locally closed and moreover it equals 
$$\cl(\tbigwedge_i a_i)\cap \op\biggl(   \tbigwedge_{x\in L} \biggl(  \big(\tbigwedge _i (b_i\to (x\vee a_i)) \big) \to \big( x\vee  \tbigwedge_j a_j\bigr) \biggr) \biggr),$$
so by assumption, it follows $$\left(\tbigwedge a_i,   \tbigwedge_{x\in L} \biggl(  \big(\tbigwedge _i (b_i\to (x\vee a_i)) \big) \to \big( x\vee  \tbigwedge_j a_j\bigr) \biggr)\right)\in U.$$

Let us now show the converse. Assume that $U$ is closed under locally exact families. We only need to show that $$\left\{ (a,b)\in \mathsf{LC}(L)\mid \cl(a)\cap\op(b)\subseteq \tbigvee_{(c,d)\in U}\cl(c)\cap\op(d)\right\}\subseteq U.$$
Let $(a,b)\in \mathsf{LC}(L)$ with $ \cl(a)\cap\op(b)\subseteq \tbigvee_{(c,d)\in U}\cl(c)\cap\op(d)\subseteq U$. 
Then
\begin{align*}\cl(a)\cap \op(b)&= \tbigvee_{(c,d)\in U}\cl(c\vee a)\cap\op(d\wedge b) \\ &=\tbigvee_{(c,d)\in U}\cl\bigl((d\wedge b)\to (c\vee a)\bigr)\cap\op\bigl((d\wedge b)\vee (d\wedge b)\to(c\vee a)\bigr)
\end{align*}
Note that $A:=\biggl\{ \bigl((d\wedge b)\to (c\vee a), (d\wedge b)\vee (d\wedge b)\to (c\vee a)\bigr) \mid (c,d)\in U\biggr\}\subseteq \cat{LC}(L).$
We first claim that $ A\subseteq U$; since $U$ is an upper set it suffices to show that for any $(c,d)\in U$, one has $(c,d)\sqsubseteq  \bigl((d\wedge b)\to (c\vee a), (d\wedge b)\vee (d\wedge b)\to (c\vee a)\bigr)$, and that is very easy to check.
Hence $A\subseteq U$ and it is a locally exact family because the corresponding join of locally closed sublocales is locally closed. Now, set
$$a_{cd}=(d\wedge b)\to (c\vee a)\quad\text{and}\quad b_{cd}= (d\wedge b)\vee (d\wedge b)\to (c\vee a),$$
$$c_0=  \tbigwedge_{(c,d)\in U} (d\wedge b)\to (c\vee a)$$
and
\begin{align*}d_0&=  \tbigwedge_{x\in L} \biggl(  \big(\tbigwedge _{(c,d)\in U}(b_{cd}\to (x\vee a_{cd})) \big) \to \big( x\vee  \tbigwedge_{(c',d')\in U} a_{c'd'}\bigr)\biggr)\\&=
  \tbigwedge_{x\in L} \biggl(  \big(\tbigwedge _{(c,d)\in U}(b\wedge d)\to (x\vee a_{cd})) \big) \to \big( x\vee  \tbigwedge_{(c',d')\in U} a_{c'd'}\bigr)\biggr)
\end{align*}
(the second equality holds by an application of \ref{H11} and \ref{H2}).  Because $U$ is closed under locally exact families, it follows that $(c_0,d_0)\in U$.

To conclude the proof it suffices to show that $(c_0,d_0)\sqsubseteq (a,b)$. First,
$$c_0=   \tbigwedge_{(c,d)\in U} b\to(d\to(c\vee a)) = b\to \bigl( \tbigwedge_{(c,d)\in U} (d\to(c\vee a))\bigr)=b\to a=a
$$
where we have used \ref{H7}, \ref{H12} and the facts that $a\in\cl(a)\cap\op(b)\subseteq \tbigvee_{(c,d)\in U} \cl(c)\cap\op(d)$ and the frame surjection corresponding to the right hand side is $\tbigwedge_{(c,d)\in U}d\to (c\vee(-))$.

Next, we have to show that $b\leq a\vee d_0$. Note that $a\leq a_{cd}$ for all $(c,d)\in U$ and so $a\leq d_0$. Hence we need to show $b\leq d_0$. Let $x\in L$, we have
\begin{align*} b\wedge  \tbigwedge _{(c,d)\in U}(b\wedge d)\to (x\vee a_{cd})&\leq b\wedge  \tbigwedge _{(c,d)\in U}(b\wedge d)\to ((b\wedge d)\to (x\vee c\vee a)) &
\\& =  b\wedge\tbigwedge _{(c,d)\in U}(b\wedge d)\to  (x\vee c\vee a) & \text{by }\ref{H7}
\\& =  b\wedge\tbigwedge _{(c,d)\in U}b\to ( d\to  (x\vee c\vee a) ) & \text{by }\ref{H7}
\\& = \tbigwedge _{(c,d)\in U}b\wedge ( d\to  (x\vee c\vee a) )& \text{by }\ref{H5}
\\& = b\wedge  \tbigwedge _{(c,d)\in U}  d\to  (x\vee c\vee a) &
\\& = b\wedge( b\to (a\vee x)) &
\\& = b\wedge (a\vee x) & \text{by }\ref{H5}
\\& = a\vee (b\wedge x)\leq(\tbigwedge_{(c',d')\in U}a_{c'd'} )\vee x.&
\end{align*}

where again we have used that the frame surjection corresponding to $\cl(a)\cap\op(b)$ is $\tbigwedge_{(c,d)\in U} d\to (c\vee (-))$.
\end{proof}

\end{document}
